\documentclass[11pt]{article}
\usepackage{amsmath,amssymb,amsthm,hyperref}
\newtheorem{lemma}{Lemma}
\newtheorem{proposition}{Proposition}
\begin{document}
\title{Erd\H{o}s Problem 868: a checked attempt}
\author{GPT\_6\_Astra\_Ultra}
\date{24 September 2026}
\maketitle

\section*{Statement and current resolution}
If an asymptotic additive basis $A\subseteq\mathbb N$ of order two has representation function tending to infinity, must it contain a minimal asymptotic basis? What if its representation function is at least $\varepsilon\log n$? Both assertions are false, by Daniel Larsen and Michael Larsen, \emph{Robust additive bases without minimal subbases}, \url{https://raw.githubusercontent.com/Larsen-Daniel/Erdos-868/main/868.pdf}. This attempt proves the obstruction mechanism and a logarithmic random-basis estimate, but does not reconstruct the probabilistic surgery needed to combine them. Its proof status is therefore unresolved, despite the known negative resolution.

We use unordered representations $r_A(n)=|\{(a,b)\in A^2:a\le b,\ a+b=n\}|$. The ordered convention in the question lies between $r_A(n)$ and $2r_A(n)$ and gives the same divergence and positive-logarithmic-lower-bound questions.

\section*{A proved obstruction criterion}
Suppose the sufficiently large integers are covered by two sets $B,C$, and $A$ is a basis with the following properties.
\begin{enumerate}
\item For each sufficiently large $c\in C$ choose at least three distinct unordered representations $c=a_i+a_i'$ in $A$. For every set $S$ formed by choosing one endpoint from all but one of these pairs, there is a witness $b=b(c,S)\in B$, with $b>2\max S$, such that the complete list of representations of $b$ in $A$ is $s+(b-s)$, $s\in S$.
\item The witnesses $b(c,S)$ tend to infinity uniformly as $c\to\infty$.
\item The smallest summand in any representation in $A$ of $b\in B$ tends to infinity as $b\to\infty$.
\end{enumerate}
Then $A$ has no minimal subbasis. To prove this, let $D\subseteq A$ be any subbasis. If a large $c\in C$ had at most one representation in $D$, from every other selected representation choose an endpoint outside $D$. This gives a set $S\subseteq A\setminus D$, and therefore a witness $b$ with no representation in $D$, a contradiction when $c$ is large. Thus $r_D(c)\ge2$ eventually on $C$. Deleting any fixed $d\in D$ destroys at most one unordered representation of any integer, so $D\setminus\{d\}$ still represents the large members of $C$. By property 3 it also still represents the large members of $B$. It is consequently a basis, and $D$ is not minimal.

For a fixed sum, distinct unordered representation pairs are disjoint as sets of summands, apart from the harmless singleton diagonal pair. Thus the endpoint selections in this criterion are well defined; they do not tacitly require independent choices among overlapping pairs. The condition $b>2\max S$ makes their newly planted witness representations distinct.

\section*{Why the logarithmic scale is compatible with a sparse basis}
Here is a complete elementary probability calculation. Select each integer $x$ independently with probability
\[
 p_x=\min\{1,K\sqrt{\log(x+1)/x}\},
\]
where $K$ is a sufficiently large constant. For a large $n$, use the pairs $\{a,n-a\}$ with $\lceil n/3\rceil\le a< n/2$. These are disjoint, so their inclusion indicators are independent. There are at least $n/7$ such pairs, and each is included with probability at least $K^2\log n/(2n)$ once $n$ is large. Their expected count is at least $K^2\log n/14$. The elementary Chernoff estimate $\mathbb P(X<\mathbb EX/2)\le\exp(-\mathbb EX/8)$ gives
\[
 \mathbb P\left(r_A(n)<K^2\log n/28\right)\le n^{-K^2/112}.
\]
Choosing, for example, $K=40$ makes these probabilities summable. The union bound on tails (equivalently the first Borel--Cantelli lemma) shows that almost surely $r_A(n)\gg\log n$ for every sufficiently large $n$. Also
\[
 \mathbb E|A\cap[1,X]|\ll_K\sqrt{X\log(X+1)}:
\]
use $p_x\le K\sqrt{\log(X+1)}x^{-1/2}$ and $\sum_{x\le X}x^{-1/2}\le2\sqrt X$. Thus a logarithmic representation lower bound is naturally available before surgery.

\section*{The precise missing step}
The random basis just constructed has not been shown to satisfy the obstruction criterion. Planting $b-S$ may introduce unwanted representations of other witnesses; deleting old representations of $b$ may destroy too many representations elsewhere. Larsen--Larsen address exactly these interactions on intervals $[X_j,X_{j+1})$, $X_j=2^{2^j}$: their Proposition 5 bounds simultaneous additive collisions, Lemma 7 bounds damage by $37$ representations, and Lemmas 9--10 provide the coverage and exact witness properties. Their construction gives the positive constant $15/(512\log2)$ in the unordered convention. These surgery estimates are reported as source results, not proved or silently assumed in this attempt.

The established progress is the complete implication from the witness mechanism to absence of minimal subbases, including deletion and quantifier checks, and the separate logarithmic random-basis estimate. The unproved compatibility of the two constructions prevents this file from claiming a complete independent counterexample.

\textbf{Completion rate estimate: 65\%.}

\end{document}
